%=============================================================================
\section{引言}
%=============================================================================

这份文档\textbf{不是}教材。教材当然好——严谨、系统，催眠效果更是一流，凌晨两点百试百灵。这是一份\textit{学习笔记}，记笔记的人学控制走的是"先把机器人搞坏，再翻书找原因"的野路子。

这份文档的目标很简单：给你够用的理论功底，让你能在真实硬件上\textbf{设计、调试、排查控制器}。数学该讲还是要讲，但每讲一个概念我们都会追问一句："机器人上场的时候，这玩意到底有什么用？"

\subsection*{这份笔记写给谁看？}

大一到大三的本科生。前置知识：学过（或正在学）微积分和基础线性代数。只要你知道导数是什么、矩阵乘法不会让你当场崩溃，就够了。

\subsection*{这份笔记的结构有什么特别之处？}

目前市面上几乎所有控制相关教材都没有类似的内容编排。

先从\textbf{数字信号处理（DSP）}讲起——控制任何东西之前，你得先能从传感器拿到干净数据。然后学如何用数学\textbf{描述系统}，说白了就是"把你的电机到底在干什么用公式写下来"。接下来是\textbf{经典控制理论}，PID 的主场——竞赛机器人的看家本领。之后是\textbf{离散化与工程实现}，在优雅的数学和跑在固定时钟频率单片机上的残酷现实之间搭一座桥。最后走进\textbf{现代控制理论}，学会那些把"能用"的机器人变成"好用"的机器人的工具：LQR、MPC 和卡尔曼滤波器。

开始吧。
